\section{CDF：把“落在左侧的概率”累计起来}
\label{sec:math-cdf}
累积分布函数（Cumulative Distribution Function, CDF）定义为
\[
 F_X(t)=P(X\leq t),\quad t\in\R.
\]
$X$是随机变量，$t$是查询位置；不要把等式右侧写成与$t$无关的固定事件。
CDF单调不减、右连续，在负无穷和正无穷处的极限分别为0和1。
它适用于离散、连续以及混合分布。

\begin{figure}[htbp]
\centering\begin{tikzpicture}
\begin{axis}[name=disc,width=72mm,height=43mm,xmin=-.3,xmax=1.3,ymin=-.05,ymax=1.12,
 xlabel={$t$},ylabel={$F_X(t)$},xtick={0,1},ytick={0,.75,1},axis lines=left]
\addplot[Teal,thick] coordinates {(-.3,0) (0,0)};
\addplot[Teal,thick] coordinates {(0,.75) (1,.75)};
\addplot[Teal,thick] coordinates {(1,1) (1.3,1)};
\addplot[only marks,mark=*,Teal,mark size=2pt] coordinates {(0,.75) (1,1)};
\addplot[only marks,mark=o,mark options={fill=white},Teal,mark size=2pt] coordinates {(0,0) (1,.75)};
\end{axis}
\begin{axis}[at={(disc.east)},anchor=west,xshift=13mm,width=72mm,height=43mm,
 xmin=-.15,xmax=.7,ymin=-.05,ymax=1.12,xlabel={$t$},xtick={0,.5},ytick={0,1},axis lines=left]
\addplot[Teal,thick] coordinates {(-.15,0) (0,0) (.5,1) (.7,1)};
\addplot[only marks,mark=*,Teal,mark size=2pt] coordinates {(0,0) (.5,1)};
\end{axis}
\end{tikzpicture}

\caption{左：$X\sim\mathrm{Bernoulli}(1/4)$的跳跃CDF。右：$U(0,0.5)$的连续CDF。实心点表示端点取值。}
\par\smallskip
\readfigure{左图在0跳到$3/4$，在1跳到1；所以$P(X=0)=3/4$。
右图在$[0,0.5]$内为$F(t)=2t$，因此$P(0.1<X\leq0.4)=0.8-0.2=0.6$。}
\end{figure}

\begin{center}
\begin{tabularx}{\textwidth}{P{30mm}Y}
\toprule
对象 & 与概率的关系\\
\midrule
离散质量函数 & $p(x)=P(X=x)$，各可能取值的质量相加为1。\\
概率密度 & 若分布具有密度$f$，则$F(t)=\int_{-\infty}^t f(u)\,du$；密度值不是单点概率。\\
累积分布函数 & 一般有$P(a<X\leq b)=F(b)-F(a)$，$a<b$。\\
\bottomrule
\end{tabularx}
\end{center}

记左极限为$F(a^-)=\lim_{t\uparrow a}F(t)$，则
\[
 P(X=a)=F(a)-F(a^-),\qquad
 P(a\leq X\leq b)=F(b)-F(a^-).
\]
例如左图$P(0<X\leq1)=1/4$，但$P(0\leq X\leq1)=1$；端点是否纳入会改变答案。
有密度的分布在单点上概率为0，此时有限个端点的开闭不改变区间概率。

\begin{pitfall}[几乎处处不等于处处]
把均匀分布密度在单个位置的值改成100，不改变任何区间积分，因而不改变分布。
若两个密度只在零测集（例如有限个点）上不同，便称为几乎处处（Almost Everywhere）相同。
类似地，ReLU只在0处不可微，不妨碍它在几乎处处可微；实现仍要约定0处如何反传。
\end{pitfall}
